\par\noindent\textbf{(1)} 推迟矢势远区与 \(\dot{\boldsymbol p}(t-r/c)\) 成正比，横向电场为 \[\boldsymbol E_{\rm rad}=\frac{1}{4\pi\varepsilon_0c^2r}\hat{\boldsymbol r}\times[\ddot{\boldsymbol p}(t-r/c)\times\hat{\boldsymbol r}]\]。对应磁场 \(\boldsymbol B_{\rm rad}=\hat{\boldsymbol r}\times\boldsymbol E_{\rm rad}/c\)。两场均横向且随 \(r^{-1}\) 衰减。
\par\noindent\textbf{(2)} 设观测极角 \(\theta\)，因偶极加速度沿 \(z\) 轴，远场幅 \(E_\theta=\omega^2p_0\sin\theta/(4\pi\varepsilon_0c^2r)\)。时间平均Poynting通量 \(\langle S_r\rangle=\varepsilon_0 cE_{\theta,0}^2/2\)，于是每立体角功率 \(\frac{dP}{d\Omega}=\frac{\omega^4p_0^2}{32\pi^2\varepsilon_0c^3}\sin^2\theta\)。
\par\noindent\textbf{(3)} 球面积分用 \(\int\sin^2\theta\,d\Omega=8\pi/3\)，得到 \(\langle P\rangle=\omega^4p_0^2/(12\pi\varepsilon_0c^3)\)。若偶极来自带电粒子位移 \(p=qx\)，则能量损失与 \(q^2\ddot x^2/(6\pi\varepsilon_0c^3)\) 的周期平均一致，辐射反作用属于 \(c^{-3}\) 阶小效应。
\par\noindent\textbf{(4)} 当 \(\theta=0,\pi\) 时横向加速度为零，故辐射沿偶极轴消失而在赤道最强。远场总功率不随 \(r\) 变，是 \(r^{-2}\) 通量和球面积 \(r^2\) 的平衡；\(\omega\to0\) 时辐射功率按 \(\omega^4\) 消失，不能用静电近场 \(r^{-3}\) 项计算远区辐射。
